\section{User Requirements Definition (Requisiti Utente \& Use Case)}

La presente sezione descrive le interazioni e le funzionalità offerte dal sistema \textbf{MyAma} dal punto di vista degli utenti. Gli Use Case sono organizzati per attore e corredati da opportuni Use Case Diagram e schede tabellari descrittive strutturate con precondizioni, scenario primario, scenari alternativi ed eccezioni, postcondizioni e requisiti associati.

% -----------------------------------------------------------
\subsection{Use Case Cittadino / Cliente}

\subsubsection{Diagramma dei Casi d'Uso: Cittadino}
\begin{figure}[H]
    \centering
    % Sostituire con l'immagine esportata da Visual Paradigm
    % \includegraphics[width=0.85\textwidth]{figure/uc_cittadino.png}
    \fbox{\parbox[c][6cm][c]{0.8\textwidth}{\centering \textsf{[ Inserire qui il Diagramma Use Case Cittadino esportato da Visual Paradigm ]}}}
    \caption{Use Case Diagram --- Cittadino / Cliente}
    \label{fig:uc_cittadino}
\end{figure}

\subsubsection{Documentazione dei Casi d'Uso: Cittadino}

\usecasetable{UC-C01: Prenota Ritiro a Domicilio}{
    Cittadino / Cliente (autenticato).
}{
    Il cittadino richiede il prelievo a domicilio di uno o più rifiuti ingombranti specificando indirizzo, tipologia del rifiuto e selezionando uno slot disponibile.
}{
    L'utente ha effettuato l'accesso al sistema.
}{
    \begin{enumerate}[leftmargin=*]
        \item Il cittadino seleziona la funzionalità di prenotazione ritiro a domicilio.
        \item Inserisce i dati del rifiuto (tipologia, descrizione, quantità, stima ingombro/peso).
        \item Inserisce l'indirizzo di ritiro e il CAP di riferimento.
        \item Il sistema valida la copertura territoriale del CAP indicato.
        \item Il sistema calcola le date e fasce orarie disponibili verificando capienza mezzi e turni autisti.
        \item Il cittadino seleziona lo slot orario desiderato.
        \item Il sistema riepiloga i dati e richiede conferma.
        \item Il cittadino conferma la prenotazione.
    \end{enumerate}
}{
    \begin{itemize}[leftmargin=*]
        \item \textbf{4a. CAP non coperto dal servizio:} Il sistema segnala l'impossibilità di procedere con il ritiro a domicilio e suggerisce il conferimento presso un'isola ecologica.
        \item \textbf{5a. Nessuno slot orario disponibile:} Il sistema invita il cittadino a selezionare un periodo alternativo.
        \item \textbf{8a. Annullamento operazione:} L'utente annulla la procedura prima della conferma; nessun dato viene salvato.
    \end{itemize}
}{
    La prenotazione viene registrata con stato \textit{Confermata} e viene generato un codice identificativo univoco inviato all'utente.
}{
    RF-01, RF-02, RF-03, RF-04, RF-05.
}

\usecasetable{UC-C02: Prenota Conferimento in Sede}{
    Cittadino / Cliente (autenticato).
}{
    Il cittadino prenota un appuntamento per conferire personalmente i rifiuti ingombranti presso un centro di raccolta AMA abilitato.
}{
    L'utente ha effettuato l'accesso al sistema.
}{
    \begin{enumerate}[leftmargin=*]
        \item Il cittadino seleziona la funzionalità di conferimento in sede.
        \item Inserisce i dati del rifiuto e indica la zona / CAP di provenienza.
        \item Il sistema mostra l'elenco delle sedi/isole ecologiche compatibili con la tipologia di rifiuto.
        \item Il cittadino seleziona la sede desiderata e visualizza gli slot temporali disponibili.
        \item Seleziona data e fascia oraria di conferimento.
        \item Conferma la richiesta di prenotazione.
    \end{enumerate}
}{
    \begin{itemize}[leftmargin=*]
        \item \textbf{3a. Nessuna sede compatibile trovata:} Il sistema notifica che per la tipologia di rifiuto indicata non vi sono centri abilitati nella zona.
        \item \textbf{5a. Slot esauriti per la sede scelta:} Il cittadino sceglie una data differente o un'altra sede.
    \end{itemize}
}{
    La prenotazione di conferimento viene registrata con stato \textit{Confermata} e viene rilasciato un pass di accesso con codice di prenotazione.
}{
    RF-06, RF-07, RF-08.
}

\usecasetable{UC-C03: Gestione e Annullamento Prenotazione}{
    Cittadino / Cliente.
}{
    Il cittadino consulta lo storico delle proprie prenotazioni attive e, se consentito dai vincoli temporali, procede alla cancellazione.
}{
    L'utente è autenticato e possiede almeno una prenotazione attiva.
}{
    \begin{enumerate}[leftmargin=*]
        \item Il cittadino accede alla sezione "Le mie prenotazioni".
        \item Seleziona una prenotazione in stato \textit{Confermata}.
        \item Richiede l'annullamento della prenotazione.
        \item Il sistema verifica che l'annullamento avvenga con almeno 24h di anticipo.
        \item Il sistema richiede conferma dell'operazione.
        \item L'utente conferma l'annullamento.
    \end{enumerate}
}{
    \begin{itemize}[leftmargin=*]
        \item \textbf{4a. Termine per l'annullamento scaduto ($<$ 24h):} Il sistema notifica che non è più possibile annullare autonomamente il servizio.
    \end{itemize}
}{
    Lo stato della prenotazione passa ad \textit{Annullata} e le risorse (slot/veicolo) vengono liberate.
}{
    RF-09, RF-10.
}

% -----------------------------------------------------------
\newpage
\subsection{Use Case Autista AMA}

\subsubsection{Diagramma dei Casi d'Uso: Autista AMA}
\begin{figure}[H]
    \centering
    % \includegraphics[width=0.85\textwidth]{figure/uc_autista.png}
    \fbox{\parbox[c][6cm][c]{0.8\textwidth}{\centering \textsf{[ Inserire qui il Diagramma Use Case Autista esportato da Visual Paradigm ]}}}
    \caption{Use Case Diagram --- Autista AMA}
    \label{fig:uc_autista}
\end{figure}

\subsubsection{Documentazione dei Casi d'Uso: Autista AMA}

\usecasetable{UC-A01: Consultazione Itinerario e Ritiri Assegnati}{
    Autista AMA.
}{
    L'autista visualizza la lista dei ritiri a domicilio programmati per il proprio turno lavorativo con dettagli su indirizzi, orari e tipologia di carico.
}{
    L'autista ha eseguito l'accesso al profilo aziendale AMA.
}{
    \begin{enumerate}[leftmargin=*]
        \item L'autista seleziona il proprio turno / data odierna.
        \item Il sistema recupera l'elenco ordinato delle prenotazioni di ritiro assegnate al suo veicolo.
        \item L'autista visualizza la scheda di dettaglio di ciascun ritiro (indirizzo, contatto, descrizione rifiuto, ingombro stimato).
    \end{enumerate}
}{
    \begin{itemize}[leftmargin=*]
        \item \textbf{2a. Nessun ritiro assegnato per la giornata:} Il sistema visualizza un messaggio informativo.
    \end{itemize}
}{
    L'autista ottiene il piano di viaggio e i dettagli operativi per lo svolgimento del servizio.
}{
    RF-11, RF-12.
}

\usecasetable{UC-A02: Registrazione Esito Ritiro}{
    Autista AMA.
}{
    Al termine dell'intervento presso il domicilio del cittadino, l'autista registra l'esito formale dell'operazione.
}{
    L'autista ha visualizzato la prenotazione assegnata ed è giunto all'indirizzo.
}{
    \begin{enumerate}[leftmargin=*]
        \item L'autista seleziona la prenotazione corrispondente.
        \item Seleziona lo stato di esito: \textit{Completato con successo}.
        \item Inserisce eventuali note operative sul carico.
        \item Invia la conferma di avvenuto ritiro.
    \end{enumerate}
}{
    \begin{itemize}[leftmargin=*]
        \item \textbf{2a. Cittadino assente o rifiuto non conforme:} L'autista registra l'esito \textit{Non Eseguito} / \textit{Respinto} specificando la motivazione formale.
    \end{itemize}
}{
    Lo stato della prenotazione viene aggiornato ad \textit{Completata} o \textit{Non Eseguita}; il cittadino riceve notifica di chiusura.
}{
    RF-13, RF-14.
}

% -----------------------------------------------------------
\newpage
\subsection{Use Case Operatore di Sede AMA}

\subsubsection{Diagramma dei Casi d'Uso: Operatore di Sede}
\begin{figure}[H]
    \centering
    % \includegraphics[width=0.85\textwidth]{figure/uc_operatore.png}
    \fbox{\parbox[c][6cm][c]{0.8\textwidth}{\centering \textsf{[ Inserire qui il Diagramma Use Case Operatore di Sede esportato da Visual Paradigm ]}}}
    \caption{Use Case Diagram --- Operatore di Sede AMA}
    \label{fig:uc_operatore}
\end{figure}

\subsubsection{Documentazione dei Casi d'Uso: Operatore di Sede}

\usecasetable{UC-O01: Accoglienza e Verifica Conferimento in Sede}{
    Operatore di Sede AMA.
}{
    L'operatore di sede verifica l'accesso del cittadino presso l'isola ecologica tramite codice prenotazione o documento e convalida il rifiuto.
}{
    Il cittadino si presenta presso il centro di raccolta con prenotazione attiva.
}{
    \begin{enumerate}[leftmargin=*]
        \item L'operatore ricerca la prenotazione tramite codice univoco o dati del cittadino.
        \item Il sistema mostra i dettagli del conferimento atteso (tipologia e quantità).
        \item L'operatore esamina fisicamente la conformità del rifiuto.
        \item Convalida la conformità e registra l'esito \textit{Completato}.
    \end{enumerate}
}{
    \begin{itemize}[leftmargin=*]
        \item \textbf{1a. Prenotazione non trovata o scaduta:} L'operatore verifica se procedere con registrazione estemporanea o rifiutare l'accesso.
        \item \textbf{3a. Rifiuto non conforme o vietato (es. tossico/amianto):} L'operatore registra la non conformità e respinge il carico.
    \end{itemize}
}{
    La prenotazione di conferimento viene chiusa con esito \textit{Completata} o \textit{Rifiutata}.
}{
    RF-15, RF-16.
}
